\chapter{Dopamine}
\index{Dopamine}

\section*{Definition and Overview}
Dopamine is a chemical messenger (neurotransmitter) in the brain and body that plays many roles, including movement, motivation, reward, pleasure, attention, and the regulation of mood. It is produced in several areas of the brain and acts on different pathways to influence behaviour and bodily functions. Disorders of dopamine are involved in many conditions, including Parkinson's disease, where dopamine-producing cells are lost, schizophrenia and psychosis, where dopamine activity is excessive, and addiction, where the reward system is hijacked by drugs. Dopamine is also used as a medicine in intensive care to support blood pressure. This chapter explains what dopamine does and the conditions related to it.

\section*{Epidemiology}
Dopamine-related conditions are common and significant. Parkinson's disease affects around 1\% of people over 60, and its motor symptoms are caused by dopamine loss. Psychotic disorders, linked to excess dopamine activity, affect around 1\% of the population. Addiction and substance use disorders affect a substantial proportion of people and involve the brain's dopamine reward system. ADHD, which involves dopamine and other transmitters, affects around 5\% of children.

\section*{Aetiology and Pathophysiology}
\begin{itemize}
    \item \textbf{Dopamine pathways:} dopamine travels along specific pathways in the brain, including those controlling movement, reward, and thinking
    \item \textbf{Parkinson's disease:} loss of dopamine-producing cells in the substantia nigra causes tremor, stiffness, and slow movement
    \item \textbf{Psychosis:} excess dopamine activity in certain pathways is linked to hallucinations and delusions
    \item \textbf{Reward and addiction:} dopamine signals pleasure and reinforces behaviours; drugs that boost dopamine drive addiction
    \item \textbf{ADHD:} altered dopamine signalling affects attention and impulse control
    \item \textbf{Medical use:} dopamine infusions raise blood pressure in critically ill patients
    \item \textbf{Other roles:} dopamine affects nausea, hormone release, and kidney blood flow
\end{itemize}

\section*{Clinical Features}
\begin{itemize}
    \item In Parkinson's disease: resting tremor, muscle stiffness, slow movement, and balance problems
    \item In psychosis: hallucinations, delusions, and disorganised thinking
    \item In addiction: compulsive drug use despite harm, cravings, and loss of control
    \item In ADHD: inattention, hyperactivity, and impulsivity
    \item As a medicine: used in hospital to support blood pressure in shock
    \item Prolactin effects: dopamine normally inhibits prolactin; disruption affects breastfeeding and periods
\end{itemize}

\section*{Differential Diagnosis}
\begin{itemize}
    \item Parkinson's disease versus other causes of tremor and movement disorders
    \item Psychosis from dopamine excess versus other causes of psychosis
    \item ADHD versus other causes of inattention in children
    \item Dopamine-related conditions versus other neurological and psychiatric disorders
\end{itemize}

\section*{When to Seek Medical Care}
Seek medical care for symptoms such as tremor, stiffness, movement difficulty, hallucinations, or significant changes in attention or behaviour. People with Parkinson's disease, psychosis, or ADHD should be under ongoing medical care.

\textbf{Contact your local emergency services for:}
\begin{itemize}
    \item Sudden weakness, facial droop, or difficulty speaking (possible stroke)
    \item Severe agitation, hallucinations, or dangerous behaviour
    \item Seizures or loss of consciousness
    \item Collapse or signs of shock
\end{itemize}

\section*{Diagnosis and Investigations}
\begin{itemize}
    \item \textbf{Clinical assessment:} dopamine-related conditions are diagnosed mainly from the history and examination
    \item \textbf{Neurological examination:} assessment of movement, tremor, and reflexes
    \item \textbf{Psychiatric assessment:} for psychosis and related symptoms
    \item \textbf{Imaging:} specialised scans such as DaTSCAN can help in Parkinson's disease
    \item \textbf{Blood tests:} prolactin levels where relevant
    \item \textbf{Specialist referral:} to neurologists and psychiatrists
\end{itemize}

\section*{Management}
\begin{itemize}
    \item \textbf{Parkinson's disease:} medicines that replace or mimic dopamine, such as levodopa, and other treatments
    \item \textbf{Psychosis:} antipsychotic medicines that block dopamine, and psychological support
    \item \textbf{ADHD:} stimulant and non-stimulant medicines that adjust dopamine and other transmitters, with behavioural support
    \item \textbf{Addiction:} treatment programs addressing the reward system, including counselling and medication
    \item \textbf{Medical use:} dopamine infusions in intensive care under specialist supervision
    \item \textbf{Exercise and therapy:} physiotherapy, occupational therapy, and psychological support
\end{itemize}

\section*{Complications}
\begin{itemize}
    \item Progressive disability in Parkinson's disease
    \item Side effects of dopamine medicines, including nausea, dyskinesias, and impulse control problems
    \item Complications of antipsychotics, including movement side effects and metabolic changes
    \item Relapse in addiction
    \item Consequences of untreated psychosis
\end{itemize}

\section*{Prognosis and Outcomes}
The prognosis depends on the condition. Parkinson's disease is progressive, but treatment controls symptoms well for many years. Psychosis and ADHD respond well to treatment, and addiction is a treatable condition with good outcomes for many people. Understanding dopamine's role has led to highly effective treatments, and ongoing research continues to improve outcomes across dopamine-related disorders.

\section*{Prevention and Self-Care}
\begin{itemize}
    \item Take prescribed medicines exactly as directed
    \item Attend specialist reviews for dopamine-related conditions
    \item Exercise regularly, which supports brain and movement health
    \item Maintain a healthy lifestyle, including sleep and stress management
    \item Avoid illicit drugs that hijack the dopamine reward system
    \item Seek help early for movement, mood, or behavioural changes
    \item Support loved ones with dopamine-related conditions and encourage treatment
\end{itemize}

\section*{Sources and Further Reading}
The following reputable sources provide further information for healthcare students and patients seeking reliable, current advice about dopamine.

\begin{itemize}
    \item Parkinson's Australia. \textit{Dopamine and Parkinson's disease.} https://www.parkinsons.org.au/
    \item SANE Australia. \textit{Psychosis and the brain.} https://www.sane.org/
    \item Healthdirect Australia. \textit{Parkinson's disease and dopamine.} https://www.healthdirect.gov.au/
    \item MSD Manual. \textit{Dopamine and brain chemistry.} https://www.msdmanuals.com/
    \item MedlinePlus. \textit{Dopamine.} https://medlineplus.gov/
    \item Alcohol and Drug Foundation. \textit{Dopamine, reward, and addiction.} https://adf.org.au/
\end{itemize}
